\lecture{1}
\title{Introduction}

\begin{document}

Class notes and course info are at
\href{https://manipulation.csail.mit.edu/}{https://manipulation.csail.mit.edu/}.
Explicitly,
\begin{itemize}
    \item This is a CI-H; writing should not be AI-generated.
    \item For homework and project implementation,
    AI usage is expected and \emph{assumed}
    (hence the 5\% problem set weighting).
\end{itemize}

We spent the rest of lecture
discussing what ``robotic manipulation'' even means;
a variety of example videos were played in lecture
that explained this well.

% There's relevantly two skills for a robot to acquire:
% dexterity (ability to perform singular tasks precisely)
% and generality (ability to perform many various tasks).

% Most (mainstream) robots are made up of the following basic components:
% \begin{itemize}
%     \item Links: pieces of connecting metal.
%     \item Joints: revoluate (rotation around a point / line)
%     and prismatic (translation along an axis).
%     \item Sensors:
%     \begin{itemize}
%         \item proploception (measures its own position, torque, etc.).
%         \item force (measures the strnegth of forces acting on itself).
%         \item tactile (measures an array of points of contact).
%         \item vision (RGB(D), Laser Scane); cameras at the wrist, e.g.
%     \end{itemize}
%     \item Actuation:
%     \begin{itemize}
%         \item Motors---very high rotational speed, very low torque.
%         \item Transmission---converting high rotational speed
%         to high torques through gear reduction, friction.
%     \end{itemize}
% \end{itemize}

\end{document}